\documentclass{article}

\usepackage{amsmath,amssymb}
\usepackage{xurl}
\usepackage[hidelinks]{hyperref}
% Shared commands for the notes in docs/paper-gaps/.
% Every note loads this file via % Shared commands for the notes in docs/paper-gaps/.
% Every note loads this file via % Shared commands for the notes in docs/paper-gaps/.
% Every note loads this file via \input{command}. Project-wide notation
% belongs here, never in per-note redefinitions.

% --- Project configuration ------------------------------------------------
\newcommand{\projectIssueBase}{https://github.com/LionSR/QICLean/issues/}
\newcommand{\projectPrBase}{https://github.com/LionSR/QICLean/pull/}
\newcommand{\projectDocBase}{https://sirui-lu.com/QICLean/docs/}
% ---------------------------------------------------------------------------

\DeclareUnicodeCharacter{2080}{\ifmmode{}_{0}\else\textsubscript{0}\fi}
\DeclareUnicodeCharacter{2081}{\ifmmode{}_{1}\else\textsubscript{1}\fi}
\DeclareUnicodeCharacter{2082}{\ifmmode{}_{2}\else\textsubscript{2}\fi}
\DeclareUnicodeCharacter{2099}{\ifmmode{}_{n}\else\textsubscript{n}\fi}

\newcommand{\Lean}{\textsc{Lean}}

% A formal declaration name, upright sans, breakable at underscores.
\ExplSyntaxOn
\NewDocumentCommand{\leanid}{m}{
  \ifmmode
    \textsf{\detokenize{#1}}
  \else
    \group_begin:
    \urlstyle{sf}
    \tl_set:Nn \l_tmpa_tl {#1}
    \tl_replace_all:Nnn \l_tmpa_tl {\_} {_}
    \exp_args:NV \path \l_tmpa_tl
    \group_end:
  \fi
}
\ExplSyntaxOff

% Traceability links; use in footnotes, not in the running prose.
\newcommand{\ghissue}[1]{\href{\projectIssueBase#1}{GitHub issue \##1}}
\newcommand{\ghpr}[1]{\href{\projectPrBase#1}{PR \##1}}
\newcommand{\doclink}[1]{\href{\projectDocBase#1.html}{\path{#1}}}

% Domain notation, using \providecommand so a note-local refinement stays
% possible.
\providecommand{\tr}{\operatorname{tr}}

% Dirac notation and the identity operator, as in blueprint/src/macros.
\usepackage{dsfont}
\providecommand{\ket}[1]{|#1\rangle}
\providecommand{\bra}[1]{\langle #1|}
\providecommand{\braket}[2]{\langle #1 | #2 \rangle}
\providecommand{\ketbra}[2]{|#1\rangle\!\langle #2|}
\providecommand{\Id}{\mathds{1}}
\providecommand{\one}{\mathds{1}}
\providecommand{\C}{\mathbb{C}}
\providecommand{\MN}[1]{M_{#1}(\C)}
\providecommand{\MD}{M_D(\C)}

% Verdict marker: \gapnote{<kind>}{<status>}. Kind is the policy
% classification (clarification, local-correction, scope-restriction,
% unfaithful, false-source, open-gap); status is open, wip (elimination
% underway), resolved, or historical. Machine-read by the site index;
% prints a verdict line.
\newcommand{\gapnote}[2]{\par\noindent\textbf{Verdict:} #1 (#2).\par}
. Project-wide notation
% belongs here, never in per-note redefinitions.

% --- Project configuration ------------------------------------------------
\newcommand{\projectIssueBase}{https://github.com/LionSR/QICLean/issues/}
\newcommand{\projectPrBase}{https://github.com/LionSR/QICLean/pull/}
\newcommand{\projectDocBase}{https://sirui-lu.com/QICLean/docs/}
% ---------------------------------------------------------------------------

\DeclareUnicodeCharacter{2080}{\ifmmode{}_{0}\else\textsubscript{0}\fi}
\DeclareUnicodeCharacter{2081}{\ifmmode{}_{1}\else\textsubscript{1}\fi}
\DeclareUnicodeCharacter{2082}{\ifmmode{}_{2}\else\textsubscript{2}\fi}
\DeclareUnicodeCharacter{2099}{\ifmmode{}_{n}\else\textsubscript{n}\fi}

\newcommand{\Lean}{\textsc{Lean}}

% A formal declaration name, upright sans, breakable at underscores.
\ExplSyntaxOn
\NewDocumentCommand{\leanid}{m}{
  \ifmmode
    \textsf{\detokenize{#1}}
  \else
    \group_begin:
    \urlstyle{sf}
    \tl_set:Nn \l_tmpa_tl {#1}
    \tl_replace_all:Nnn \l_tmpa_tl {\_} {_}
    \exp_args:NV \path \l_tmpa_tl
    \group_end:
  \fi
}
\ExplSyntaxOff

% Traceability links; use in footnotes, not in the running prose.
\newcommand{\ghissue}[1]{\href{\projectIssueBase#1}{GitHub issue \##1}}
\newcommand{\ghpr}[1]{\href{\projectPrBase#1}{PR \##1}}
\newcommand{\doclink}[1]{\href{\projectDocBase#1.html}{\path{#1}}}

% Domain notation, using \providecommand so a note-local refinement stays
% possible.
\providecommand{\tr}{\operatorname{tr}}

% Dirac notation and the identity operator, as in blueprint/src/macros.
\usepackage{dsfont}
\providecommand{\ket}[1]{|#1\rangle}
\providecommand{\bra}[1]{\langle #1|}
\providecommand{\braket}[2]{\langle #1 | #2 \rangle}
\providecommand{\ketbra}[2]{|#1\rangle\!\langle #2|}
\providecommand{\Id}{\mathds{1}}
\providecommand{\one}{\mathds{1}}
\providecommand{\C}{\mathbb{C}}
\providecommand{\MN}[1]{M_{#1}(\C)}
\providecommand{\MD}{M_D(\C)}

% Verdict marker: \gapnote{<kind>}{<status>}. Kind is the policy
% classification (clarification, local-correction, scope-restriction,
% unfaithful, false-source, open-gap); status is open, wip (elimination
% underway), resolved, or historical. Machine-read by the site index;
% prints a verdict line.
\newcommand{\gapnote}[2]{\par\noindent\textbf{Verdict:} #1 (#2).\par}
. Project-wide notation
% belongs here, never in per-note redefinitions.

% --- Project configuration ------------------------------------------------
\newcommand{\projectIssueBase}{https://github.com/LionSR/QICLean/issues/}
\newcommand{\projectPrBase}{https://github.com/LionSR/QICLean/pull/}
\newcommand{\projectDocBase}{https://sirui-lu.com/QICLean/docs/}
% ---------------------------------------------------------------------------

\DeclareUnicodeCharacter{2080}{\ifmmode{}_{0}\else\textsubscript{0}\fi}
\DeclareUnicodeCharacter{2081}{\ifmmode{}_{1}\else\textsubscript{1}\fi}
\DeclareUnicodeCharacter{2082}{\ifmmode{}_{2}\else\textsubscript{2}\fi}
\DeclareUnicodeCharacter{2099}{\ifmmode{}_{n}\else\textsubscript{n}\fi}

\newcommand{\Lean}{\textsc{Lean}}

% A formal declaration name, upright sans, breakable at underscores.
\ExplSyntaxOn
\NewDocumentCommand{\leanid}{m}{
  \ifmmode
    \textsf{\detokenize{#1}}
  \else
    \group_begin:
    \urlstyle{sf}
    \tl_set:Nn \l_tmpa_tl {#1}
    \tl_replace_all:Nnn \l_tmpa_tl {\_} {_}
    \exp_args:NV \path \l_tmpa_tl
    \group_end:
  \fi
}
\ExplSyntaxOff

% Traceability links; use in footnotes, not in the running prose.
\newcommand{\ghissue}[1]{\href{\projectIssueBase#1}{GitHub issue \##1}}
\newcommand{\ghpr}[1]{\href{\projectPrBase#1}{PR \##1}}
\newcommand{\doclink}[1]{\href{\projectDocBase#1.html}{\path{#1}}}

% Domain notation, using \providecommand so a note-local refinement stays
% possible.
\providecommand{\tr}{\operatorname{tr}}

% Dirac notation and the identity operator, as in blueprint/src/macros.
\usepackage{dsfont}
\providecommand{\ket}[1]{|#1\rangle}
\providecommand{\bra}[1]{\langle #1|}
\providecommand{\braket}[2]{\langle #1 | #2 \rangle}
\providecommand{\ketbra}[2]{|#1\rangle\!\langle #2|}
\providecommand{\Id}{\mathds{1}}
\providecommand{\one}{\mathds{1}}
\providecommand{\C}{\mathbb{C}}
\providecommand{\MN}[1]{M_{#1}(\C)}
\providecommand{\MD}{M_D(\C)}

% Verdict marker: \gapnote{<kind>}{<status>}. Kind is the policy
% classification (clarification, local-correction, scope-restriction,
% unfaithful, false-source, open-gap); status is open, wip (elimination
% underway), resolved, or historical. Machine-read by the site index;
% prints a verdict line.
\newcommand{\gapnote}[2]{\par\noindent\textbf{Verdict:} #1 (#2).\par}


\title{The Trimming Radius in the Bond-Trimming Lemma\\
    of Landau, Vazirani and Vidick}
\author{}
\date{2026-10-25}

\begin{document}
\maketitle
\gapnote{local-correction}{resolved}

\section{The source statement}

For a vector $u$ with Schmidt decomposition $u=\sum_j\lambda_j a_j\otimes b_j$
across a cut, $\lambda_1\ge\lambda_2\ge\cdots$, the trimmed vector
$\mathrm{trim}_D(u)=\sum_{j\le D}\lambda_j a_j\otimes b_j$ keeps the $D$
largest Schmidt terms. The bond-trimming lemma of Landau, Vazirani and
Vidick reads as follows.\footnote{The source is
    \cite{Landau2013PolynomialTime}. Only version~1 of arXiv:1307.5143 is
    listed, and its source file \path{algo53.tex} is the one checked here:
    trimming is defined at lines~330--333, the Eckart--Young lemma
    \path{lem:eckart} is at lines~335--339, and the lemma \path{l:trim1}
    with its proof is at lines~531--542. The published versions (ITCS 2014,
    doi 10.1145/2554797.2554825; Nature Physics 11, 566 (2015), doi
    10.1038/nphys3345) could not be checked.}
Given a vector $v$ with $D$ nonzero Schmidt coefficients across the cut, for
every $u$,
\begin{align}
    \bigl|\langle\mathrm{trim}_{D/\varepsilon}(u),v\rangle\bigr|
     & \ge|\langle u,v\rangle|-\varepsilon.
    \label{eq:lvv13_trimming_radius_printed}
\end{align}
The vectors $u$ and $v$ are unit vectors in every application.

The proof argues by contradiction. Put $w=u-\mathrm{trim}_{D/\varepsilon}(u)$
and assume $|\langle w,v\rangle|>\varepsilon$. The Schmidt coefficients of $w$
are $\lambda_{\lceil D/\varepsilon\rceil+1},\lambda_{\lceil
D/\varepsilon\rceil+2},\dots$, so the Eckart--Young theorem bounds the overlap
of $w$ with the unit vector $v$ of Schmidt rank $D$ by
\begin{align}
    |\langle w,v\rangle|
     & \le\Bigl(\sum_{i=1}^{D}\lambda_{\lceil D/\varepsilon\rceil+i}^2\Bigr)^{1/2}.
    \label{eq:lvv13_trimming_radius_eckart}
\end{align}
The assumption therefore gives a tail sum larger than $\varepsilon^2$. The
source instead records it as larger than $\varepsilon$, and only that stronger
claim makes the closing estimate
\begin{align}
    \sum_{j=1}^{\lceil D/\varepsilon\rceil}\lambda_j^2
     & \ge\Bigl\lceil\frac1\varepsilon\Bigr\rceil
    \sum_{i=1}^{D}\lambda_{\lceil D/\varepsilon\rceil+i}^2>1
    \label{eq:lvv13_trimming_radius_closing}
\end{align}
contradict the normalization of $u$. With the correct tail bound
$\varepsilon^2$, the right-hand side is only larger than $\varepsilon$, and
no contradiction follows.

\section{The counterexample}

The statement (\ref{eq:lvv13_trimming_radius_printed}) itself fails, not
only its proof. Take $\varepsilon=1/k$ with an integer $k\ge2$, so that
$D/\varepsilon=kD$ is an integer, and let $u$ have $N=(k+1)D$ equal Schmidt
coefficients,
\begin{align}
    u & =\frac{1}{\sqrt N}\sum_{j=1}^{N}a_j\otimes b_j,
      & v & =\frac{1}{\sqrt D}\sum_{j=kD+1}^{(k+1)D}a_j\otimes b_j.
    \label{eq:lvv13_trimming_radius_counterexample}
\end{align}
Both are unit vectors, and $v$ has exactly $D$ nonzero Schmidt coefficients.
Trimming $u$ to $kD$ terms removes precisely the terms in $v$, so
\begin{align}
    \langle\mathrm{trim}_{kD}(u),v\rangle & =0,
                                           &
    \langle u,v\rangle                     & =\sqrt{\frac DN}
    =\sqrt{\frac{\varepsilon}{1+\varepsilon}}.
    \label{eq:lvv13_trimming_radius_overlaps}
\end{align}
The printed bound (\ref{eq:lvv13_trimming_radius_printed}) would force
$\sqrt{\varepsilon/(1+\varepsilon)}\le\varepsilon$, that is
$\varepsilon^2+\varepsilon\ge1$, which fails for every
$\varepsilon<(\sqrt5-1)/2$, in particular for every $\varepsilon=1/k$ with
$k\ge2$. The smallest instance is $D=1$, $\varepsilon=1/2$: the vector
$u=(a_1\otimes b_1+a_2\otimes b_2+a_3\otimes b_3)/\sqrt3$ is trimmed to two
terms, $v=a_3\otimes b_3$, the trimmed overlap is $0$, and
$|\langle u,v\rangle|-\varepsilon=1/\sqrt3-1/2>0$.
This is a genuine error in the source rather than a degenerate reading of
it: the data are unit vectors with nonzero Schmidt coefficients, and the
failure persists for every $D$.

\section{The correction}

Trimming to $k$ terms costs at most
\begin{align}
    \bigl|\langle u-\mathrm{trim}_k(u),v\rangle\bigr|
     & \le\|u\|\,\|v\|\sqrt{\frac{D}{k+1}}
    \label{eq:lvv13_trimming_radius_corrected}
\end{align}
for every $u$, every $v$ of Schmidt rank at most $D$, and every $k$; the
vectors need not be normalized. Consequently,
$|\langle\mathrm{trim}_k(u),v\rangle|\ge|\langle u,v\rangle|-
\|u\|\|v\|\sqrt{D/(k+1)}$. For unit vectors, the radius
$k+1\ge D/\varepsilon^2$ gives error at most $\varepsilon$. The
counterexample (\ref{eq:lvv13_trimming_radius_counterexample}) shows that a
radius linear in $1/\varepsilon$ does not suffice, since the loss there is
$\sqrt{D/(k+D)}$ at radius $k$.

The proof runs as follows. Let $\rho$ be the reduced density matrix of $u$ on
the left factor, with eigenvalues $\mu_1\ge\mu_2\ge\cdots$, and let $P$ be the
projection that performs the trimming, so that $\operatorname{tr}(P\rho)$ is
the sum of the $k$ largest eigenvalues. Let $T$ be the projection onto the
range of $(I-P)$ applied to the coefficient matrix of $v$; it has rank at
most $D$ and is orthogonal to $P$. By the Cauchy--Schwarz inequality
$|\langle u-\mathrm{trim}_k(u),v\rangle|^2\le\operatorname{tr}(T\rho)\,
\|v\|^2$. Ky Fan's maximum principle applied to the projection $P+T$ of rank
at most $k+D$ gives
\begin{align}
    \operatorname{tr}(T\rho)
     & \le\sum_{i=k+1}^{k+D}\mu_i\le D\mu_{k+1}
    \le\frac{D}{k+1}\operatorname{tr}\rho
    =\frac{D}{k+1}\|u\|^2.
    \label{eq:lvv13_trimming_radius_ky_fan}
\end{align}
The argument does not depend on which eigenvectors of a repeated eigenvalue
the trimming keeps, consistent with the footnote at line~332 of the source.

The later trimming lemma of the same authors with Arad has the quadratic
scaling of (\ref{eq:lvv13_trimming_radius_corrected}).\footnote{See
    \cite{Arad2016RigorousRG}, version~2, source file
    \path{local_alg_arxiv_v2.tex}, Lemma \path{lem:trimming},
    lines~660--697.}
There, a threshold $\xi$ on the Schmidt coefficients gives Schmidt rank at
most $s\xi^{-2}$ and error $\sqrt{\ell bs}\,\xi$, so the rank grows as the
inverse square of the error.

\section{Effect downstream}

Claim \path{c:trim} of the source (lines~549--567) applies the lemma at each
of the $n$ bonds with $D=r(n)$ and $\varepsilon=1/(48n)$, so with radius
$p_2(n)=48n\,r(n)$, and sums the $n$ errors to $n/(48n)=1/48$. With the
corrected radius $D/\varepsilon^2$ the same argument needs
$p_2(n)=(48n)^2r(n)$. This is still a polynomial, so Claim \path{c:trim}
and the polynomial running time of the algorithm survive with the larger
$p_2$. The companion lemma \path{l:trim2} (lines~544--547), that trimming
across one cut does not increase the Schmidt rank across any other cut, is
unaffected.

\section{Formalization}

The trimming projection of a vector at $D$ is
\leanid{Matrix.IsTrimmingProjection}: an orthogonal projection of rank at
most $D$ on the left factor whose trace against the reduced density matrix
is the sum of its $D$ largest eigenvalues. Every choice in a degenerate
Schmidt coefficient is admitted, and \leanid{Matrix.exists_isTrimmingProjection}
gives existence for every $D$. The trimmed vector is
\leanid{Matrix.leftApply}. The Eckart--Young lemma is
\leanid{Matrix.IsTrimmingProjection.norm_inner_le_re_inner_normalized},
with the real overlap computed by
\leanid{Matrix.IsTrimmingProjection.inner_normalized_eq_sqrt}. The corrected
bound (\ref{eq:lvv13_trimming_radius_corrected}) is
\leanid{Matrix.IsTrimmingProjection.norm_inner_sub_leftApply_le}, and its
form as a lower bound on the trimmed overlap is
\leanid{Matrix.IsTrimmingProjection.norm_inner_sub_le_norm_inner_leftApply}.
The printed statement (\ref{eq:lvv13_trimming_radius_printed}) is not
formalized.

\bibliographystyle{alpha}
\bibliography{references}

\end{document}
